% GENERATED by docs/paper/make_tex_tables_revision.py -- do not edit by hand.
\begin{table}[htbp]
\centering
\caption{Per-crop zero-shot accuracy on MobileCLIP2-S0/dfndr2b over a common pool of 78 classes (chance 1.3\%).}
\label{tab:loco}
\begin{tabular}{lrrl}
\toprule
Crop & $N$ & Zero-shot $\uparrow$ & 95\% CI \\
\midrule
Orange & 1,361 & 28.9\% & [26.4\%, 31.3\%] \\
Coffee & 1,582 & 15.8\% & [14.0\%, 17.7\%] \\
Apple & 10,000 & 12.7\% & [12.0\%, 13.3\%] \\
Peach & 1,107 & 12.3\% & [10.4\%, 14.2\%] \\
Potato & 2,698 & 9.5\% & [8.3\%, 10.7\%] \\
Corn & 9,403 & 4.8\% & [4.4\%, 5.3\%] \\
\midrule
\textbf{Pooled} & \textbf{26,151} & \textbf{10.5\%} & [10.2\%, 10.9\%] \\
\bottomrule
\end{tabular}
\begin{flushleft}\footnotesize No crop is held out and no model is refit here: the descriptor head is training-free, so this is one zero-shot run broken down by crop. \emph{Trained} labels crops that sit in the seen pool of the main experiments. The two hardest crops are both \emph{trained} crops, so difficulty does not follow the held-out boundary. The held-out crops here are configuration A's three, so this bounds that configuration and not B or C; they nevertheless average above the trained pool here (19.0\% against 9.0\%), so this split is not adversarially hard. This run covers six of the study's 18 crops, so the held-versus-trained contrast rests on three crops a side. Measured with the \texttt{rich} strategy, whose prototypes collide (Section~\ref{sec:strategies}), so the held/trained gap is itself subject to that confound.\end{flushleft}
\end{table}
